\clearpage
\refstepcounter{sourcechapter}
\phantomsection
\section*{评论}
\label{sec:comments}
\addcontentsline{toc}{section}{评论}
\MBSetSourceChapterAnchor{comments}

\hypertarget{comments-page-396}{}
\subsection*{第 1 章}

第 \ref{sec:ch1-function-spaces} 节对基本空间 $V$ 和 $H$ 作了预备研究. 迹定理由 J.L. Lions and E. Magenes 的方法证明, 见文献 [1]. 这里给出的 $H^1$ 的刻画基于 G. de Rham 的流理论中的一个定理. O.A. Ladyzhenskaya [1]对 $n=3$ 给出了更初等的证明. R. Temam [9]给出了适用于所有维数的 O.A. Ladyzhenskaya 证明的简化版本. 注 1.9 给出了另一种避开 de Rham 定理的方法; 另见命题 1.1 之前脚注的末尾.

本书没有系统研究或综述 Sobolev 空间, 只在需要时回顾这些空间的性质, 特别是在第 \ref{sec:ch1-function-spaces} 节和第 \ref{sec:ch2-discrete-inequalities-compactness} 节. 如正文所述, 关于证明和进一步材料, 读者可参阅 R.S. Adams [1], S. Agmon [1], J.L. Lions [1], J.L. Lions and E. Magenes [1], J. Nečas [1], L. Sobolev [1]等.

Stokes 方程的变分形式最早由 J. Leray [1, 2, 3]在非线性情形的一般框架中引入, 用以研究 Navier--Stokes 方程的弱解或湍流解. Stokes 变分问题解的存在性很容易由经典投影定理得到; 为完整起见, 本书回顾了该定理的证明. 非变分 Stokes 问题及解的正则性研究基于 L. Cattabriga [1]关于 $n=3$ 的论文, 以及 S. Agmon, A. Douglis and L. Nirenberg [1]关于任意维椭圆系统的论文; 本书回顾了这些结果而未给出证明. 关于正则性的另一种方法, 见 V.A. Solonnikov and V.E. Scadilov [1]. 另见 V.A. Solonnikov [4], I.I. Vorovich and V.I. Yudovich [1].

赋范空间和变分问题的逼近概念尤其由 J.P. Aubin [1]和 J. Cea [1]研究; 本书采用 R. Temam [8]的表述. 离散 Poincaré 不等式以及用有限差分逼近 $V$ 的内容见 J. Cea [1]. 用协调有限元逼近 $V$ 最早由 M. Fortin [2]研究并使用; 本书对协调有限元逼近 (APX2), (APX3)的描述基本遵循 M. Fortin [2]. 该文献还包含用这种离散化作计算的许多结果. 使用泡函数的思想来自 P.A. Raviart; 这里给出的逼近 (APX2')的表述是新的. 逼近 (APX4)由 J.P. Thomasset [1]研究并使用. 关于以非协调有限元逼近无散向量函数的材料来自 M. Crouzeix, R. Glowinski, P.A. Raviart 和作者.\footnote{致读者: 这些评论来自本书初版(1977). 较新的评论见原书第 337 页和附录 III.}

\hypertarget{comments-page-397}{}
关于该主题的其他方面, 包括高次非协调有限元和更精细的误差估计, 可见 M. Crouzeix and P.A. Raviart [1]. 关于数值实验, 见 F. Thomasset [2]; 关于轴对称三维流的情形, 另见 P. Lailly [1].

关于有限元在流体力学中的其他应用, 见 J.T. Oden, O.C. Zienkiewicz, R.H. Gallagher and T.D. Taylor [1], 以及 1976 年 6 月在意大利举行的会议论文集(待出版). 关于有限元一般理论, 可提及 I. Babuska and A.K. Aziz [1], P.G. Ciarlet [1], P.A. Raviart [2], G. Strang and G. Fix [1]的综合性著作, 以及 A.K. Aziz [1]主编的论文集. 关于一般情形下有限元的更多文献, 读者可参阅这些著作的参考文献. 这里对有限元方法的叙述几乎完全自足: 只假设了少数特殊结果, 因为证明它们需要引入与本书范围相距甚远的工具.

离散 Stokes 问题后, 需要求解一个有限维线性问题, 其中未知量 $u_h$ 属于有限维空间 $V_h$. 有两种可能:
\begin{enumerate}
\item[(a)] 空间 $V_h$ 具有自然且简单的基, 使问题化为关于 $u_h$ 在该基下各分量的稀疏矩阵线性系统; 此时通过求解该线性系统来求解问题.
\item[(b)] 若非如此, 即使有限维问题具有唯一解, 求解也不简单, 因为矩阵可能病态或不稀疏. 此时必须引入适当算法求解这些问题; 这正是第 \ref{sec:ch1-numerical-algorithms} 节的目的.
\end{enumerate}

第 \ref{sec:ch1-numerical-algorithms} 节描述的算法由 K.J. Arrow, L. Hurwicz and H. Uzawa [1]在优化理论和经济学框架中引入. 这些过程在流体动力学问题中的应用见 J. Céa, R. Glowinski and J.C. Nedelec [1], M. Fortin [2], M. Fortin, R. Peyret and R. Temam [1]. 关于参数 $\varrho$(或 $\varrho$ 与 $\alpha$)最优选择的实验研究, 见 D. Bégis [1], M. Fortin [2]; Crouzeix [2]在一个很特殊的情形中给出了该问题的理论解答.

用罚方法逼近不可压缩流体最早由 R. Temam [2a, 2b]研究. 这里给出的 $u_\epsilon$ 的完整渐近展开来自 M.C. Pelissier [1].

\subsection*{第 2 章}

第 \ref{sec:ch2-existence-uniqueness} 节给出若干关于非线性稳态 Navier--Stokes 方程解存在唯一性的标准结果. 我们主要遵循 O.A. Ladyzhenskaya [1]和 J.L. Lions [2]. 关于解正则性和流体动力学位势理论的更完整讨论见 O.A. Ladyzhenskaya [1]; 关于正则性另见 H. Fujita [1]. 无界区域中的稳态 Navier--Stokes 方程由 R. Finn [1]--[5], R. Finn and D.R. Smith [1, 2], J.G. Heywood [1, 3]研究.

关于稳态 Navier--Stokes 方程的一些较新理论结果见 C. Foias and R. Temam [2, 3], C. Foias and J.C. Saut [2], J.C. Saut and R. Temam [2], D. Serre [1, 2, 3], R. Temam [11, 16].

\hypertarget{comments-page-398}{}
第 \ref{sec:ch2-discrete-inequalities-compactness} 节给出离散 Sobolev 不等式和紧性定理, 其证明技术性很强. 有限差分情形的证明原则与连续情形的相应证明平行, 见例如 J.L. Lions [1], J.L. Lions--E. Magenes [1]. 离散 Sobolev 不等式的证明此前未曾发表; 离散紧性定理的证明见 P.A. Raviart [1]. 对协调有限元, 证明简单得多: 特别地, 对离散紧性定理, 一个简单技巧便把问题化为连续情形. 对非协调有限元, Sobolev 不等式的证明基于非协调有限元理论的特殊技术. 离散紧性定理通过比较协调元和非协调元证明; 这些结果是新的.

稳态 Navier--Stokes 方程离散化的讨论遵循第 \ref{chap:steady-state-stokes} 章发展出的原则. 一般收敛定理与第 \ref{chap:steady-state-stokes} 章类似, 并考虑同样类型的 $V$ 的离散化; 差别来自精确问题的解可能不唯一. 第 \ref{subsec:ch2-stationary-numerical-algorithms} 节的数值算法由 M. Fortin, R. Peyret and R. Temam [1]引入并检验. 三线性形式 $b$ 的修改见 \eqref{eq:ch2-antisymmetrized-form}, 它对应于引入稳定项 $\frac12(\operatorname{div}u)u$, 以及当函数不是无散函数时引入其离散类似物; 这一修改由 R. Temam [2a, 2b, 3, 4]引入并使用.

近年来人们研究了 Navier--Stokes 方程及相关方程稳态解的非唯一性. 这一方向的主要结果来自 P.H. Rabinowitz [2]和 W. Velte [1, 2]. 在 [2]中, Rabinowitz 通过显式构造两个不同解, 证明对流问题解的非唯一性; 第一个是流体静止时的平凡解, 第二个由迭代过程构造. W. Velte 的工作基于拓扑方法, 分岔理论和拓扑度理论. [1]中考虑的问题与 P.H. Rabinowitz [2]中的对流问题相同. 在 [2]中, W. Velte 证明 Taylor 问题解的非唯一性; 该情形与第 \ref{sec:ch2-existence-uniqueness} 节证明存在性的那个问题非常相似, 但并不相同. 第 \ref{sec:ch2-bifurcation-nonuniqueness} 节紧密遵循这一表述. 关于分岔理论的其他应用, 特别参见 J.B Keller and S. Antman [1], L. Nirenberg [1], P.H. Rabinowitz [4, 6], 以及 \emph{Rocky Mountain Journal of Mathematics} (1973)第 3 卷第 2 期.

\subsection*{第 3 章}

线性化 Navier--Stokes 方程的存在唯一性结果, 即第 \ref{sec:ch3-linear-case} 节, 是线性变分方程解的一般存在唯一性结果的特例, 见例如 J.L. Lions--E. Magenes [1, vol. 2]. 为完整起见, 本书对若干技术结果给出了初等证明; 它们通常是更深结果的简单推论. 例如引理 1.1 在向量值分布理论 L. Schwartz [2]的框架中更自然, 引理 1.2 可由插值方法 J.L. Lions--E. Magenes [1]证明.

定理 2.1 是非线性演化方程理论中使用的标准紧性定理之一. 其他紧性定理在 J.L. Lions [2]中得到证明和应用. 这些结果的一种较新推广见 R. Temam [16].

\hypertarget{comments-page-399}{}
第 \ref{sec:ch3-existence-uniqueness} 节和第 \ref{sec:ch3-semidiscrete-existence} 节给出的非线性 Navier--Stokes 方程存在唯一性结果现已成为经典, 它们延续了 J. Leray [1, 2, 3]的早期工作; 见 E. Höpf [1, 2], O.A. Ladyzhenskaya [1], J.L. Lions [2, 3], J.L. Lions and G. Prodi [1], J. Serrin [3]. 关于解正则性的进一步结果, 以及 Navier--Stokes 方程经典可微解的存在性研究, 可见 O.A. Ladyzhenskaya [1]第 2 版. 关于解的解析性, 见 C. Foias and G. Prodi [1], H. Fujita and K. Masuda [1], C. Kahane [1], K. Masuda [1], J. Serrin [3], C. Foias and R. Temam [4].

还应提及本书没有讨论的两种完全不同的存在唯一性方法. 第一种是 E.B. Fabes, B.F. Jones and N.M. Riviere [1]的方法, 它基于奇异积分算子并在 $L^p$ 空间中给出存在唯一性结果. 另一种是 V. Arnold [1]以及 D.G. Ebin and J. Marsden [1]的方法, 它把 Navier--Stokes 初值问题同 Riemann 流形上的测地线相联系, 从而使用整体分析方法.

第 \ref{sec:ch3-general-stability-convergence} 节关于 Navier--Stokes 方程简单离散格式稳定性和收敛性的讨论基本是新的. R. Temam [2a, 2b, 3, 4]对不同方程或不同格式作了类似研究. O.A. Ladyzhenskaya [5]给出若干无条件稳定单步格式的稳定性与收敛性; 关于分步格式另见 A.J. Chorin [2], O.A. Ladyzhenskaya and V.I. Rivkind [1]. 除 A.J. Chorin [2]外, 所有这些文献都像本书一样, 通过获得近似解的适当先验估计并使用紧性定理来证明收敛性. 在 [2]中, A.J. Chorin 假设存在非常光滑的解, 再比较近似解与精确解.

第 \ref{sec:ch3-projection-method} 节的开头部分实质上是后续两步投影格式的引言. 第 \ref{subsec:ch3-projection-two-step} 节描述的分步格式, 即投影方法, 由 A.J. Chorin [1, 2, 3]和作者 R. Temam [3]独立引入. A.J. Chorin 考虑的格式略有不同, 没有稳定项 $\frac12(\operatorname{div}u)u$, 即没有以 $\widehat b$ 代替 $b$. 该格式的应用及其他方面特别见 C.K. Chu and G. Johansson [1], C.K. Chu, K.V. Morton and K.V. Roberts [1], M. Fortin, R. Peyret and R. Temam [1], M. Fortin [1], M. Fortin and R. Temam [1], G. Marshall [1, 2], C.S. Peskin [1]. 这种格式推广了 G.I. Marchuk [1]和 N.N. Yanenko [1]引入并研究的分步法, 见第 \ref{sec:ch3-artificial-compressibility} 节.

用微可压缩流体方程逼近 Navier--Stokes 方程, 即第 \ref{subsec:ch3-ac-perturbed-problems} 节, 由 A.J. Chorin [1]和 R. Temam [3]独立引入. N.N. Yanenko [1]考虑了略微更复杂的摄动方程. 引入这些摄动后可以使用分步法, 其研究见第 \ref{subsec:ch3-ac-limit} 节. 应指出, 第 \ref{sec:ch3-projection-method} 节的格式本身就是分步格式, 不需要考虑摄动方程.

这里给出的分步格式收敛性证明来自 R. Temam [3, 4], 并遵循 R. Temam [1]引入的方法. 关于分步法的其他方面, 见 G.I. Marchuk [1], N.N. Yanenko [1, 2]及其参考文献; 另见 R. Temam [1, 6, 7]. 其他类型的摄动问题并非为了应用分步法, 而是为了克服约束 $\operatorname{div}u=0$ 造成的困难; 见 J.L. Lions [4], R. Temam [2a, 2b]. 关于交替方向方法和分步法的进一步结果, 见 O.A. Ladyzhenskaya and V.I. Rivkind [1], V.I. Rivkind and B.S. Epstein [1], B.S. Epstein [1].

\hypertarget{comments-page-400}{}
第 \ref{sec:ch3-general-stability-convergence} 节至第 \ref{sec:ch3-artificial-compressibility} 节的材料只是流体力学方程逼近方面大量工作的很小一部分. O.M. Belotserkovskii [1]主编的论文集, M. Holt [2], H. Cabannes and R. Temam [1], R.D. Richtmyer [1], F. Thomasset [1], T. Kawai [1]中可见当时最新的结果和很有用的文献. 另见 Los Alamos Scientific Laboratory 汇编的参考文献表.

这里使用的方法还可处理许多其他问题. 对 Navier--Stokes 方程本身, 可考虑不同边界条件, 见 Iooss [1]; 周期解, 见 G. Prouse [1, 2]; 以及变分不等式, 见 J.L. Lions [2]. 随机 Navier--Stokes 方程见 A. Bensoussan and R. Temam [1], C. Foias [1], C. Foias and R. Temam [5, 9], M.I. Vishik and A.V. Fursikov [1, 2, 3]. 受 Navier--Stokes 方程支配的系统的最优控制问题见 M. Cuvelier [1]; 关于较新结果, 见第 \ref{sec:app3-other-problems} 节末尾.

Navier--Stokes 方程数学理论遇到的困难, 使若干作者重新审视导出这些方程的流体力学假设, 并提出具有较好数学性质的新模型; 见 S. Kaniel [1], O.A. Ladyzhenskaya [1].

涉及其他方程的类似模型通常把 Navier--Stokes 方程与其他方程耦合. 其中包括对流方程, 其处理与 Navier--Stokes 方程几乎相同; 多种流体模型; 污染模型 G. Marshall [1]或血流模型 C.S. Peskin [1]; 以及具有浓度方程形式的海洋学模型. 更复杂的模型是磁流体力学方程和 Bingham 方程, 见 G. Duvaut and J.L. Lions [1, 2]; 后者是非 Newton 流体的一个例子.

本书没有展开 Euler 方程的数学理论. 关于基于解析方法的处理, 见 C. Bardos [1], T. Kato [1, 2], J.L. Lions [2], R. Temam [10, 12], V.I. Yudovich [1].

关于 $\nu\to0$ 时 Navier--Stokes 方程行为的一些结果见 J.L. Lions [2], V.I. Yudovich [1]. 与 Burgers 方程有关的一个模型方程的类似问题, 已由 C.M. Brauner, P. Penel and R. Temam [1], P. Penel [1]完整研究; 另见 R. Temam [12]中 C. Bardos, U. Frish, P. Penel and P.L. Sulem 的工作.

\hypertarget{comments-page-401}{}
\clearpage
\thispagestyle{empty}
\null

\hypertarget{comments-page-402}{}
\clearpage
\section*{第 3 版(修订版)补充评论}
\label{sec:additional-comments-third-edition}
\addcontentsline{toc}{subsection}{第 3 版(修订版)补充评论}

这里给出 Navier--Stokes 方程理论和数值分析中若干最新结果的说明. 这些结果主要集中在三个方向.

\paragraph{(a) 解的存在性, 唯一性和正则性.}
自 J. Leray [1, 2, 3]和 E. Hopf [1]的工作以来, 已知对含时 Navier--Stokes 方程, 若数据足够光滑, 则初值问题存在唯一光滑解, 它定义在某个时间区间 $(0,T^*)$ 上; 此后该解可延拓为正则性可能较低的解, 见第 \ref{sec:ch3-existence-uniqueness} 节和第 \ref{sec:ch3-semidiscrete-existence} 节. 目前仍不知道解是否对所有时间保持光滑. 沿 B. Mandelbrot [1, 2]的思想, 近来人们研究了解奇异集, 即速度为无穷的点集, 的 Hausdorff 维数; 见 V. Scheffer [1]--[4], C. Foias and R. Temam [4], 以及 L. Caffarelli, R. Kohn and L. Nirenberg [1]的最新论文, 后者包含当时关于奇异集 Hausdorff 维数的最佳估计.

关于解存在性和正则性的其他较新结果包括:
\begin{itemize}
\item 有界区域中稳态流解集的研究, 见 C. Foias and J.C. Saut [2], C. Foias and R. Temam [2, 3], J.C. Saut and R. Temam [2].
\item 对应于非光滑数据, 特别是非光滑区域, 的解的存在性与正则性. 这适用于 Couette--Taylor 流或腔体流等经典情形, 见 D. Serre [2, 3]. 关于无界区域中的流, 还应提及 D. Serre [1]的结果: 在某些情形中, 他找到了解空间中的整条直线, 这对非退化非线性问题很不寻常.
\item 含时 Navier--Stokes 方程弱解的一些新先验估计, 它们蕴含 $L^\infty$ 范数关于时间属于 $L^1$, 即在三维空间中 $u\in L^1(0,T;L^\infty(\Omega)^3)$; 见 C. Foias, C. Guillopé and R. Temam [1].
\item 相容性条件的推导. 这些条件是数据必须满足的充要条件, 从而保证含时方程的解在 $t=0$ 附近具有正则性; 这当然与 $t>0$ 时可能出现的奇性无关. 见 R. Temam [15].
\end{itemize}

\paragraph{(b) 长时间行为与湍流.}
若体积力与时间无关, 则时间不在 Navier--Stokes 方程中显式出现, 方程成为自治的无穷维动力系统. 与理解湍流现象有关的一个问题, 是含时 Navier--Stokes 方程解在 $t\to\infty$ 时的行为.

\hypertarget{comments-page-403}{}
近来人们研究了 Navier--Stokes 方程的渐近分析: 各种范数在无穷远的界, 流动决定模(或决定参数)的数目, 吸引子的结构和性质等. 见 A.V. Babin and M.I. Vishik [1]--[4], P. Constantin, C. Foias, O. Manley and R. Temam [1], P. Constantin, C. Foias and R. Temam [2], C. Foias and R. Temam [4, 10], C. Foias and J.C. Saut [1], C. Guillopé [1], C. Foias, O. Manley, R. Temam and Y. Trève [1], E. Lieb [1], D. Ruelle [1], R. Temam [16], O.A. Ladyzhenskaya [6, 7], I.M. Vishik [2].

\paragraph{(c) 数值逼近.}
已经出现了大量关于 Navier--Stokes 方程数值逼近的论文. 它们尤其研究有限元方法, 有限元实现的实际问题, 罚方法在流动问题中的应用, 见第 \ref{sec:ch1-penalty-method} 节, 以及 Galerkin 逼近解在长时间区间上的行为. 在众多参考文献中, 可见 M. Bercovier [1], V. Girault and P.A. Raviart [1], R. Glowinski [1], F. Thomasset [2], T. Kawai [1], J.G. Heywood and R. Rannacher [1], P. Constantin, C. Foias and R. Temam [1], 以及这些文献所列参考文献. 讨论计算流体力学其他方面的专著包括 M. Holt [4], D. Gottlieb and S. Orszag [1], R. Peyret and T.D. Taylor [1]; 另见这些文献的参考文献.
